\begin{scanpage}{0205}{187}
The secondary edge wave propagated from the end $z=z_2$ is excited during diffraction at this end of the primary current wave
\[-S\psi_-^0(z-z_1)e^{iw_0z_1+ik(z-z_1)},\tag{30.06}\]
where by $\psi_\pm^0(z)$ we mean the functions obtained from the functions $\psi_\pm(z)$ by replacing $\vartheta$ by $\vartheta_0$. For the purpose of calculating the desired secondary waves, it is necessary for us, first of all, to find that external field which, when applied to an infinite conductor $(-\infty<z<\infty)$, would excite the current (30.06) on its section $(z_1<z<\infty)$.

For this purpose, let us study the current induced in an infinite conductor by the external field
\[
E_z^e=\hat E_{0z}e^{iw_0z}\cdot\varepsilon(z-z_1),\quad
\varepsilon(z-z_1)=\begin{cases}1&\text{with }z<z_1,\\0&\text{with }z>z_1.\end{cases}\tag{30.07}
\]
Let us assume that $w_0$ has a small negative imaginary part $(\operatorname{Im}w_0\leq0)$. We may regard the quantity $\hat E_{0z}e^{iw_0\zeta}d\zeta$ as a concentrated emf which, in accordance with equation (28.03), creates in an infinite conductor
\[
\frac{c\hat E_{0z}}{4\ln\dfrac{i}{\gamma ka}}\psi(|z-\zeta|)e^{iw_0\zeta+ik|z-\zeta|}\,d\zeta.\tag{30.08}
\]
Therefore, in accordance with the principle of superposition, the total current created in the region $z_1<z<\infty$ by the external field (30.07) will equal
\[
\begin{aligned}
J(z)&=\frac{c\hat E_{0z}}{4\ln\dfrac{i}{\gamma ka}}\int_{-\infty}^{z_1}\psi(z-\zeta)e^{iw_0\zeta+ik(z-\zeta)}\,d\zeta\\
&=\frac{c\hat E_{0z}}{4\ln\dfrac{i}{\gamma ka}}e^{iw_0z}\cdot\int_{z-z_1}^{\infty}\psi(\zeta)e^{i(k-w_0)\zeta}\,d\zeta.
\end{aligned}\tag{30.09}
\]
The resulting integral may be expressed in terms of the functions $\psi(z-z_1)$ and $\psi_-^0(z-z_1)$, and the corresponding relationships derived in \S\ 2 of [82]. As a result, we find
\end{scanpage}
